\documentclass[11pt,a4paper]{article}
\usepackage{fontspec,amsmath,amssymb,amsthm,microtype}
\setmainfont{lmroman10-regular.otf}[BoldFont=lmroman10-bold.otf,ItalicFont=lmroman10-italic.otf,BoldItalicFont=lmroman10-bolditalic.otf]
\usepackage[margin=25mm]{geometry}
\usepackage[unicode,hidelinks]{hyperref}
\hypersetup{pdfauthor={LangevinLean Project}}
\setlength{\emergencystretch}{3em}
\setlength{\parskip}{0.4em}
\setlength{\parindent}{0pt}
\allowdisplaybreaks[2]
\raggedbottom
\title{\texorpdfstring{The Gaussian weak equation for bounded tests with finite gradient energy}{The Gaussian weak equation for bounded tests with finite gradient energy}}
\author{LangevinLean Project}
\date{}
\begin{document}\maketitle
Draft mathematical exposition; these notes may contain mistakes.


Write \(g(x)=\nabla U(x)\), \(c_t(x)=x-tg(x)\), and \(K_t(x,y)=N_t e^{-\|y-c_t(x)\|^2/(4t)}\). For fixed starting point, the spatial flux is \(F_{t,x}=\nabla_yK_t+K_tg(x)\). Direct coordinate differentiation and an orthonormal-basis sum prove \(\operatorname{div}_yF_{t,x}=\dot K_t\). This agrees with diffusion coefficient one and drift \(-g(x)\).



Let \(\psi\) be bounded and \(C^1\), with \(\nabla\psi\in L^2(\rho_t)\), and assume finite initial gradient energy. Gaussian second moments and Cauchy–Schwarz prove integrability of the flux-test product, the flux-gradient product, and the flux derivative times the test. Global coordinate integration by parts is therefore valid. Summing coordinates and using Fubini yields \(\int\dot p_t\psi\,dy=-\iint\langle F_{t,x},\nabla\psi\rangle d\mu\,dy\). The equation follows from the explicit Gaussian kernel; no Hessian of the potential or assumed weak equation is used.



Earlier results used here: \cite[Section~3.38]{LangevinLean2026ULASC}.

\section{\texorpdfstring{Choose the spatial flux with the correct Fokker–Planck sign}{Choose the spatial flux with the correct Fokker–Planck sign}}


The component density evolves under frozen drift \(-g(x)\). Its time derivative is \(\Delta_yK_t+\operatorname{div}_y(K_tg(x))\), so define

\[F_{t,x}(y)=\nabla_yK_t(x,y)+K_t(x,y)g(x)=K_t(x,y)\left(g(x)-\frac{y-c_t(x)}{2t}\right).\]

The last expression defines this flux. Its divergence equals the density's time derivative, while the weak equation consequently carries a minus sign after integration by parts. The frozen gradient \(g(x)\) is constant in the spatial variable \(y\), which avoids differentiating the possibly merely Lipschitz gradient field.



The flux uses the frozen starting gradient \(\nabla U(x)\), not the current gradient at the observation. Its divergence equals the kernel time derivative, while the difference from the target drift later appears as entropy forcing.


\section{\texorpdfstring{Differentiate each coordinate of the flux explicitly}{Differentiate each coordinate of the flux explicitly}}


Let \(r=y-c_t(x)\). For a direction \(v\), differentiating the kernel and the affine residual gives

\[D_v\langle F_{t,x}(y),v\rangle=K_t(x,y)\left(\frac{\langle r,v\rangle^2}{4t^2}-\frac{\langle r,v\rangle\langle g(x),v\rangle}{2t}-\frac{\|v\|^2}{2t}\right).\]

The affine vector factor has derivative \(-v/(2t)\) in direction \(v\). Pairing with the fixed \(v\) contributes \(-\|v\|^2/(2t)\). The product rule contributes the derivative of the kernel times \(\langle g-r/(2t),v\rangle\). Inserting \(D_vK_t=-K_t\langle r,v\rangle/(2t)\) and expanding gives the other two terms. The frozen \(g(x)\) is constant in \(y\), so no second derivative of the potential is taken.



Differentiate each flux component using the Gaussian gradient and product rules. The resulting diagonal quadratic displacement and drift terms are the coordinate contributions to its divergence.


\section{\texorpdfstring{Sum the directional derivatives to obtain a coordinate-\allowbreak{}free divergence}{Sum the directional derivatives to obtain a coordinate-free divergence}}


Sum over an orthonormal basis \((b_i)_{i=1}^d\). Parseval gives \(\sum_i\langle r,b_i\rangle^2=\|r\|^2\), completeness gives \(\sum_i\langle r,b_i\rangle\langle g,b_i\rangle=\langle r,g\rangle\), and \(\sum_i\|b_i\|^2=d\). Thus

\[\operatorname{div}F_{t,x}=K_t\left(-\frac d{2t}+\frac{\|r\|^2}{4t^2}-\frac{\langle r,g\rangle}{2t}\right)=K_t\left(-\frac d{2t}+\frac{\|y-x\|^2}{4t^2}-\frac{\|g\|^2}{4}\right).\]

For the last equality substitute \(r=y-x+tg\) and expand its squared norm. Its mixed term cancels the displayed pairing with \(g\), leaving the squared-position and squared-gradient terms. Positive time justifies the denominator operations. The resulting coefficient is the same one obtained by differentiating the kernel in time.



Summing over an orthonormal coordinate basis collects displacement squares into a Euclidean squared norm and diagonal constants into the dimension. This gives a coordinate-free divergence formula.


\section{\texorpdfstring{The density of the initial/\allowbreak{}output pair}{The density of the initial/output pair}}


The pair \((X,Y_t)\), constructed from independent initial state and standard noise, satisfies

\[\mathcal L(X,Y_t)(dx,dy)=K_t(x,y)\,\mu(dx)\,dy.\]

The identity for the conditional Gaussian starting-proposal pair-density formula already proves this transport for a measurable center map. Apply it to \(x\mapsto x-t\nabla U(x)\), whose measurability follows from the Lipschitz gradient. It then checks that the real kernel here is exactly the previously normalized proposal density. This identity transfers integrability proved on the independent initial/\allowbreak{}noise probability to the kernel-weighted integrals used in integration by parts.



The independent state-noise pair is pushed to the starting-proposal pair. Its conditional Gaussian density is the kernel, so nonnegative and integrable pair expectations have the stated Lebesgue representation.



Earlier results used here: \cite[Section~3.34.1]{LangevinLean2026ULASC}; \cite[Section~3.37.12]{LangevinLean2026ULASC}; \cite[Section~3.37.13]{LangevinLean2026ULASC}; \cite[Section~3.38.1]{LangevinLean2026ULASC}.

\section{\texorpdfstring{Prove flux pairings integrable by transporting two L\ensuremath{{}^{2}} factors}{Prove flux pairings integrable by transporting two L² factors}}


For a continuous scalar field \(\varphi\in L^2(\rho_t)\), the coordinate flux pairing is integrable:

\[\int K_t(x,y)\left|\left\langle g(x)-\frac{y-c_t(x)}{2t},v\right\rangle\varphi(y)\right|\,d\mu(x)\,dy<\infty.\]

Under the independent pair \((X,Z)\), the vector factor is \(g(X)-(\sqrt{2t}/(2t))Z\). Both terms are in \(L^2\), by finite initial gradient energy and the Gaussian second moment. Its pairing with fixed \(v\) is therefore in \(L^2\). Pushforward change of variables gives \(\varphi(Y_t)\in L^2\). Cauchy–Schwarz proves integrability of their scalar product. The pair-density identity gives the displayed weighted Lebesgue integral, without a pointwise decay assumption on the test.



The flux pairing can be transported to a product of square-integrable fields on the independent-noise source. Cauchy–Schwarz proves absolute integrability there and hence in the kernel representation.



Earlier results used here: \cite[Section~3.34.1]{LangevinLean2026ULASC}; \cite[Section~3.36.5]{LangevinLean2026ULASC}; \cite[Section~3.37.4]{LangevinLean2026ULASC}.

\section{\texorpdfstring{Admit the gradient pairing of a C\ensuremath{{}^{1}} finite-\allowbreak{}energy test}{Admit the gradient pairing of a C¹ finite-energy test}}


Choose the scalar field \(\varphi(y)=D\psi(y)[v]=\langle\nabla\psi(y),v\rangle\). For \(\psi\in C^1\), this field is continuous. If \(\nabla\psi\in L^2(\rho_t)\), its coordinate pairing is also in \(L^2\).



For a fixed direction \(v\), the derivative of the \(C^1\) test is \(D\psi(y)[v]=\langle\nabla\psi(y),v\rangle\), a continuous function. Its absolute value is bounded by \(\|v\|\|\nabla\psi(y)\|\), so it has finite square energy. The flux-pairing estimate then gives the required integrability of \(F_i\partial_i\psi\).



A bounded \(C^1\) test with finite gradient square energy lies in the intended weak-test class. The preceding flux estimate verifies integrability of its gradient pairing before integration by parts.


\section{\texorpdfstring{Control the boundary product from boundedness of the test}{Control the boundary product from boundedness of the test}}


A bounded continuous test satisfies \(|\psi|\le M\), hence belongs to \(L^2(\rho_t)\) because \(\rho_t\) is a probability measure. The general flux-pairing result therefore gives integrability of \(\langle F_{t,x}(y),v\rangle\psi(y)\).



A bounded test has finite square energy under the interpolation probability. Pair it with the finite-energy flux by Cauchy–Schwarz to obtain integrability of \(F_i\psi\), the boundary product required for global integration by parts. No separate pointwise decay of the test at infinity is assumed.



The test itself is bounded, so multiplying by an integrable flux component preserves integrability. This supplies the boundary-product condition in the global coordinate integration-by-parts theorem.


\section{\texorpdfstring{Match the divergence to the independently differentiated Gaussian kernel}{Match the divergence to the independently differentiated Gaussian kernel}}


Direct differentiation of the normalized Gaussian kernel gives its time derivative:

\[\dot K_t(x,y)=K_t(x,y)\left(-\frac d{2t}+\frac{\|y-x\|^2}{4t^2}-\frac{\|\nabla U(x)\|^2}{4}\right).\]

Summing the coordinate flux derivatives gives the same scalar expression as the kernel time derivative. The identity follows by expanding the centered displacement square and canceling its mixed terms. Thus the weak PDE uses a proved equality of explicit derivatives, rather than an assumed evolution equation.



The divergence and time derivative have been calculated independently from their explicit formulas. Algebraic comparison identifies them pointwise, fixing both the diffusion and drift signs in the weak equation.


\section{\texorpdfstring{Dominate each unit-\allowbreak{}coordinate flux derivative by a quadratic Gaussian energy}{Dominate each unit-coordinate flux derivative by a quadratic Gaussian energy}}


For \(\|v\|=1\), Cauchy–Schwarz bounds \(|\langle r,v\rangle|\le\|r\|\) and \(|\langle g,v\rangle|\le\|g\|\). The inequality \(\|r\|\|g\|\le\|r\|^2+\|g\|^2\) then gives

\[\left|D_v\langle F_{t,x},v\rangle\right|\le\left(\frac1{4t^2}+\frac1{2t}\right)K_t(x,y)(1+\|g(x)\|^2+\|r\|^2).\]

Apply the triangle inequality to the three coordinate derivative terms. The squared pairing is bounded by \(\|r\|^2\), and the mixed pairing is bounded by \(\|r\|\|g\|\le\|r\|^2+\|g\|^2\). Each resulting scalar coefficient is bounded by the sum \(1/(4t^2)+1/(2t)\). Multiplying by \(K_t\ge0\) gives the displayed inequality. Only quadratic gradient energy is required.



A derivative of one flux coordinate is bounded by the kernel times a constant plus quadratic displacement and gradient energies. These are precisely the terms included in the established Gaussian energy envelope.



Earlier results used here: \cite[Section~3.37.4]{LangevinLean2026ULASC}.

\section{\texorpdfstring{Integrability of the coordinate flux derivatives}{Integrability of the coordinate flux derivatives}}


The common envelope \(K_t(1+\|g(X)\|^2+\|Y-c_t(X)\|^2)\) is jointly integrable: its integral is \(1+G_0+2td\), where \(G_0=\int\|\nabla U\|^2d\mu\). The preceding pointwise bound therefore proves joint integrability of every unit-coordinate flux derivative.



The explicit polynomial-Gaussian derivative formula is jointly measurable because the gradient and the elementary kernel operations are measurable. Its absolute value is bounded by a finite time coefficient times the Gaussian weight just integrated. Domination therefore proves absolute integrability of the coordinate derivative, as required before Fubini and integration by parts.



Integrating the quadratic envelope first in the Gaussian observation and then in the initial law gives a finite value from dimension and \(G_0\). This verifies absolute integrability of each coordinate derivative.


\section{\texorpdfstring{Apply global integration by parts after proving all three product integrabilities}{Apply global integration by parts after proving all three product integrabilities}}


For one unit coordinate \(v\), global integration by parts gives

\[\int D_v\langle F_{t,x}(y),v\rangle\,\psi(y)\,d\mu(x)\,dy=-\int\langle F_{t,x}(y),v\rangle D_v\psi(y)\,d\mu(x)\,dy.\]

The coordinate derivative times the bounded test is integrable by derivative domination. The flux times the directional test derivative is integrable by the finite-gradient-energy pairing estimate, and the flux times the test is integrable by boundedness. Fubini provides these three section conditions for almost every initial state. At such a state both functions are differentiable, so global directional integration by parts applies. Integrating its section identity over the initial law gives the displayed equality. The required condition at infinity is integrability of these products; no additional pointwise test decay is assumed.



The product, its directional derivative, and the flux-gradient pairing are all integrable. These checks make global integration by parts applicable without assuming compact support of the bounded test.


\section{\texorpdfstring{The weak equation for bounded continuously differentiable tests}{The weak equation for bounded continuously differentiable tests}}


Sum the coordinate identities over an orthonormal basis. The divergence identity turns the sum on the left into \(\dot K_t\), while completeness of the basis turns the sum on the right into \(\langle F_{t,x},\nabla\psi\rangle\). Fubini then yields

\[\int\dot p_t(y)\psi(y)\,dy=-\int\langle F_{t,x}(y),\nabla\psi(y)\rangle\,d\mu(x)\,dy.\]

Each coordinate integrand is integrable by the preceding estimates, so finite sums commute with their integrals. The time derivative of the kernel is integrable, and multiplication by the bounded test preserves that property. Fubini consequently exchanges the starting-state and observation integrals and identifies their mixture with \(\dot p_t\). All integrability conditions follow from the Gaussian law, initial gradient energy, boundedness of the test, and finite test-gradient energy.



Sum the coordinate identities and use divergence equal to the kernel time derivative. Fubini identifies the mixture derivative on the left, producing the weak equation for every admitted bounded finite-energy test.


\begin{thebibliography}{ULA-LSI}
\bibitem[ULA-SC]{LangevinLean2026ULASC}
LangevinLean Project. \emph{Unadjusted Langevin under Strong Convexity}. LL-01, 2026. Draft report, content revision 6f6f50decdb6. \url{https://odakyildiz.github.io/LangevinLean/docs/pdf/UnadjustedLangevin/FullProof.pdf}.
\end{thebibliography}

\end{document}
